\begin{proposition}[Exclusion of primary Kodaira general fibres]\hypertarget{tag-01F5}{}\label{cdp:st:C.3.10}
\StatementRecord{cdp:st:C.3.10}{proved}{01F5}{\SeriesRef{cdp:st:C.3.9}, \SeriesRef{cdp:st:D.1.12}, \SeriesRef{cdp:st:D.1.9}, \SeriesRef{cdp:st:D.1.7}, \SeriesRef{cdp:eq:C.PK.4}, \SeriesRef{cdp:st:D.2.3}, \SeriesRef{cdp:eq:C.PK.5}, \SeriesRef{cdp:eq:C.PK.1}, \SeriesRef{cdp:eq:C.PK.2}, \SeriesRef{cdp:eq:C.PK.7}; \cite[pp.~373--374]{Borcea1984}; \cite[\S1]{SS2002}; \cite[Remark~1]{Borcea1984}; \cite[\S3]{Borcea1984}; \cite[Theorem~2.18 and \S\S2.19--2.21]{Steenbrink1976}}


A smooth compact complex threefold with $b_1(Y)=b_2(Y)=0$ and
$a(Y)=1$ cannot have a holomorphic algebraic reduction
$f:Y\to\PP^1$ whose general fibre is a primary Kodaira
surface. Together with \ref{cdp:st:C.3.9} this excludes both Kodaira
types using only rational Betti-number vanishing.
\end{proposition}


\begin{proof}

\begin{enumerate}[leftmargin=*]
\item Let $B=\PP^1\setminus\Sigma$ be the smooth base and write
$P_c=f^{-1}(c)$ and $\pi_c:P_c\to E_c$ for the canonical
elliptic algebraic reduction of a smooth fibre. The degree
$d>0$ is constant. All the smooth fibres have the fixed
primary Kodaira differentiable type; every complex structure
on that type is again a primary Kodaira structure
\cite[pp.~373--374]{Borcea1984}.
Use the intrinsic quotient and its orientation from \eqref{cdp:eq:C.SK.2}:
\[
 Q_c=\Pic^0(P_c)/\pi_c^*\Pic^0(E_c),\qquad
 q_c:Q_c\xrightarrow{\simeq}\C^*.
\]
The class of the pullback of a degree-$d$ line has coordinate
$\alpha(c)\in\Delta^*$. The holomorphic moduli comparison
in the proof of \ref{cdp:st:C.3.9} applies directly here, so this bounded
function extends to $\PP^1$ and has a single nonzero
constant value $\alpha$.

\item Put $\mathscr E=R^1f_*\OO_Y$ and
$\mathscr G=R^2f_*\OO_Y$. The primary Kodaira case of
\ref{cdp:st:D.1.12} proves, on the original threefold and every full fibre,
\begin{equation}
 \operatorname{rank}\mathscr E=2,\qquad
 \operatorname{rank}\mathscr G=1,\qquad
 \bigwedge^2\mathscr E\xrightarrow{\simeq}\mathscr G.
\hypertarget{tag-00V8}{}\label{cdp:eq:C.PK.1}
\end{equation}
Both sheaves are locally free and satisfy base change.
Let $\gamma$ be the integral transgression section of \ref{cdp:st:D.1.9},
and let $\beta\in H^0(\PP^1,\mathscr E)=H^1(Y,\OO_Y)$
be its coefficient image. This section is nowhere zero.
For $u\in H^1(P_c,\OO_{P_c})$ define the complex-linear
functional $\lambda_c$ by the exact identity
\begin{equation}
 q_c\bigl([\exp(2\pi\mathrm i z u)]\bigr)
     =\exp(2\pi\mathrm i z\lambda_c(u))
     \quad(z\in\C).
\hypertarget{tag-00V9}{}\label{cdp:eq:C.PK.2}
\end{equation}
Its kernel is the line
$A_c=\pi_c^*H^1(E_c,\OO_{E_c})$.
The function $k(c)=\lambda_c(\beta_c)$ takes integral
values, because $\beta_c$ is the coefficient image of
an integral class. It is holomorphic, by the actual relative
Picard quotient in \eqref{cdp:eq:C.SK.2}, hence is a constant integer $k$.

\item When $k\ne0$, choose $n>0$ with $nd\ge2$ and choose
$z\in\C$ with $\exp(2\pi\mathrm i kz)=\alpha^n$.
The ambient line $M=\exp(2\pi\mathrm i z\beta)$ then
has, on every smooth fibre, the same $Q_c$ class as the
pullback of a degree-$nd$ line on $E_c$. Its difference
from that pullback belongs to $\pi_c^*\Pic^0(E_c)$.
Thus
$M|_{P_c}=\pi_c^*L_c$ for a line $L_c$ of degree $nd$,
and $h^0(P_c,M|_{P_c})=nd\ge2$. The relative linear
system gives $a(Y)\ge2$, as at the end of \ref{cdp:st:C.3.9}.
\item For $k=0$, $\beta_c\in A_c$. The coefficient map
$H^1(E_c;\R)\to H^1(E_c,\OO_{E_c})$ is an isomorphism
of real vector spaces. Its pullback, together with the real
injectivity of \ref{cdp:st:D.1.7} on $P_c$, proves the exact implication
\begin{equation}
 \beta_c\in A_c
 \quad\Longrightarrow\quad
 \gamma_c\in\pi_c^*H^1(E_c;\R).
\hypertarget{tag-00VA}{}\label{cdp:eq:C.PK.3}
\end{equation}
The lattice $\pi_c^*H^1(E_c;\Z)$ is saturated in
$H^1(P_c;\Z)$: the elliptic bundle induces a surjection
on fundamental groups, so the cokernel of the induced
injection on integral $H^1$ is torsion free. Thus $\gamma$
lies in this rank-two integral local system. Divide it by
its fixed positive divisibility to obtain a primitive fixed
section $\gamma_0$. On a simply connected open set choose
a complementary integral section $\delta$ with oriented
intersection number one. The ratio of their nonzero
coefficient images is a holomorphic elliptic period
\[
 \tau=\frac{\iota(\delta)}{\iota(\gamma_0)},
 \qquad \operatorname{Im}\tau>0,
\]
after choosing the sign convention of the ordered pair.
Monodromy fixes $\gamma_0$ and replaces $\delta$ by
$\delta+N\gamma_0$, $N\in\Z$. Therefore
$\exp(2\pi\mathrm i\tau)$ is a single-valued bounded
holomorphic function on $B$. It extends to $\PP^1$ and
is constant. Its logarithmic derivative gives $d\tau=0$,
and a nonzero integral translation cannot fix this constant
period. Consequently the Albanese integral $H^1$ local
system has identity monodromy, and $j(E_c)$ is constant.

\item Both moduli $(j,\alpha)$ and the degree $d$ are now
constant. Borcea's classification, in the precise form
\cite[\S1]{SS2002}, makes all smooth fibres biholomorphic
to one primary Kodaira surface $P$. Borcea's complete
holomorphic marked period family \cite[Remark~1]{Borcea1984}
expresses our family, near any smooth base point, as the
pullback by a holomorphic map into his period domain
$\mathcal D\subset\PP(H^2(P;\C))$. By his Theorem~3,
constant isomorphism class means that the image of this
map lies in one orbit of the integral cohomology action.
That orbit is countable, since the action factors through
$\operatorname{GL}(H^2(P;\Z)/\operatorname{torsion})$.
A connected subset of a countable subset of a complex
manifold is one point: a real coordinate separating two
points would have a nontrivial interval as its connected
image. Therefore the period parameter is constant and
its pullback is a product. 
\item Write the fundamental group as
\begin{equation}
 \Gamma=\langle x,y,w,u\mid
       w,u\text{ central},\ [x,y]=w^d\rangle.
\hypertarget{tag-00VB}{}\label{cdp:eq:C.PK.4}
\end{equation}
Its centre $\langle w,u\rangle$ is the fundamental group
of the elliptic fibre, and its quotient is the fundamental
group of the elliptic base. A biholomorphism inducing the
identity on the base $H^1$ induces the identity on this
quotient. The relation in \eqref{cdp:eq:C.PK.4} then forces it to fix
$w$. On an elliptic fibre its linear part is a holomorphic
group automorphism. Such an automorphism acts by a complex
scalar on the universal cover; fixing the nonzero period
$w$ forces that scalar to be one. It therefore fixes both
$w$ and $u$. The remaining possible actions on $x,y$ are
exactly multiplication by central elements.

\item In the real
nilmanifold presentation choose invariant one-forms
$a,b,c,e$ with
\[
 da=db=de=0,\qquad dc=d\,a\wedge b.
\]
The four classes $[a\wedge c],[b\wedge c],
[a\wedge e],[b\wedge e]$ form a basis of
$H^2(P;\R)$, as in \cite[\S3]{Borcea1984}.
An automorphism of \eqref{cdp:eq:C.PK.4} with the stated action extends
to the real nilpotent group: in logarithmic coordinates it
fixes the two central generators and adds central vectors
to the two base generators. Its pullback fixes $a,b$ and
replaces each of $c,e$ by itself plus a real linear
combination of $a,b$. Each of the four displayed classes
therefore changes only by a multiple of the exact form
$a\wedge b$. Since the universal cover of $P$ is
contractible, a map of $P$ and this group-model map with
the same induced fundamental-group map are homotopic.
This computes the action of the original biholomorphism
and proves it is the identity on $H^2$.
In particular every monodromy of our isotrivial family,
whose action on the base $H^1$ was just proved to be the
identity, acts identically on $H^2(P;\C)$.

\item Resolve all exceptional
fibres by successive blowups with smooth centres contained
in those fibres, writing
$\rho:\widetilde Y\to Y$ and $h=f\rho$. For each such
blowup, $\rho_*\OO_{\widetilde Y}=\OO_Y$ and
$R^j\rho_*\OO_{\widetilde Y}=0$ for $j>0$.
One local verification uses the formal neighbourhoods of
the exceptional projective bundle: their successive
conormal layers restrict to nonnegative twists
$\OO_{\PP^{r-1}}(l)$, $l\ge0$, which have no higher
cohomology. The analytic theorem on formal functions
gives the asserted vanishing; normality gives the
degree-zero equality. The Leray edge is consequently the
canonical isomorphism
\begin{equation}
 R^2h_*\OO_{\widetilde Y}\xrightarrow{\simeq}
 R^2f_*\OO_Y=\mathscr G.
\hypertarget{tag-00VC}{}\label{cdp:eq:C.PK.5}
\end{equation}
The same identities hold for a composite by derived
pushforward.

\item Let $D=(h^{-1}(\Sigma))_{\mathrm{red}}$ and use the
relative logarithmic de Rham complex
\[
 \mathscr K^\bullet=
 \Omega_{\widetilde Y}^\bullet(\log D)\big/
 \bigl(h^*\Omega_{\PP^1}^1(\log\Sigma)
       \wedge\Omega_{\widetilde Y}^{\bullet-1}(\log D)\bigr).
\]
Projection onto degree zero is a morphism of complexes
$\mathscr K^\bullet\to\OO_{\widetilde Y}[0]$:
it is the identity in degree zero and zero in every
positive degree. The exact analytic comparison of \ref{cdp:st:D.2.3},
namely \cite[Theorem~2.18 and \S\S2.19--2.21]{Steenbrink1976},
applies to this resolution with all its fibre multiplicities
and without a K\"ahler hypothesis. It gives the logarithmic
de Rham bundle
$\mathscr H^2=\mathbb R^2h_*\mathscr K^\bullet$.
Identity monodromy makes every residue zero, by the
residue interval $[0,1)$ and the exponential comparison.
The global flat frame on $B$ therefore extends across
every point of $\Sigma$, proving
$\mathscr H^2\simeq H^2(P;\C)\otimes\OO_{\PP^1}$.
Taking degree-two direct hypercohomology of the displayed
projection and then \eqref{cdp:eq:C.PK.5} gives a holomorphic map
\begin{equation}
 p:H^2(P;\C)\otimes\OO_{\PP^1}
          \longrightarrow\mathscr G.
\hypertarget{tag-00VD}{}\label{cdp:eq:C.PK.6}
\end{equation}
Over $B$ it is precisely the coefficient projection
$H^2(P_c;\C)\to H^2(P_c,\OO_{P_c})$.

\item In an isotrivial chart this latter map is constant. It is
surjective: the class of
$\overline\alpha\wedge\overline\eta$ used in the proof
of \ref{cdp:st:D.1.12} is represented by a closed invariant two-form
and has nonzero Dolbeault class. The charts also show
that its kernel is a flat subspace. Choose a constant flat
class $v_0\in H^2(P;\C)$ whose projection is nonzero,
and put $s=p(v_0)\in H^0(\PP^1,\mathscr G)$.
This section is nowhere zero on $B$.

\item By \eqref{cdp:eq:C.PK.1} and the nowhere-zero section $\beta$, the
following is an exact sequence of vector bundles on the
entire projective line:
\begin{equation}
 0\longrightarrow\OO_{\PP^1}\beta
  \longrightarrow\mathscr E
  \xrightarrow{\ \beta\wedge-\ }\mathscr G
  \longrightarrow0.
\hypertarget{tag-00VE}{}\label{cdp:eq:C.PK.7}
\end{equation}
Since $H^1(\PP^1,\OO)=0$, choose
$v\in H^0(\PP^1,\mathscr E)=H^1(Y,\OO_Y)$ with
$\beta\wedge v=s$.
We identify exactly the scalar that this lift induces in
the quotient \eqref{cdp:eq:C.PK.2}. In each isotrivial chart,
$\beta$ is the coefficient image of a fixed integral
class, $v_0$ is a fixed de Rham class, and $s=p(v_0)$
is a fixed nonzero class in $H^2(P,\OO_P)$.
The intrinsically oriented quotient $q_c$ is also fixed
by the chart: every automorphism preserves its
positive-degree class $\alpha\in\Delta^*$, so it cannot
induce inversion of $\C^*$.
Both linear maps $u\mapsto\beta\wedge u$ and
$u\mapsto\lambda(u)s$ have kernel $A$ and are nonzero.
Thus they differ by a constant $\kappa\in\C^*$ on
that chart. The constants agree on overlaps, because
all the maps and sections used in their definition are
the actual global ones. The connectedness of $B$ gives
one constant $\kappa$, and \eqref{cdp:eq:C.PK.7} yields
\begin{equation}
 \lambda_c(v_c)=\kappa^{-1}\ne0
 \quad\hbox{for every }c\in B.
\hypertarget{tag-00VF}{}\label{cdp:eq:C.PK.8}
\end{equation}

\item Choose $n>0$ with $nd\ge2$ and $z\in\C$ with
$\exp(2\pi\mathrm i z\kappa^{-1})=\alpha^n$.
The single ambient line $M=\exp(2\pi\mathrm i z v)$
has this quotient coordinate on every smooth fibre.
Exactly as in the $k\ne0$ case,
$M|_{P_c}=\pi_c^*L_c$ with $\deg L_c=nd\ge2$.
Its direct image has rank at least two and its relative
linear system gives a surface image over $\PP^1$,
contradicting $a(Y)=1$.

\end{enumerate}
\end{proof}